\chapter{Introduction}
In modern times the world moves forward with incredible speed and creates new buildings aimed to engage people who use them, but often these environments are designed for adults with fully developed psychological and physical capabilities.
People who don't have the necessary requirements to navigate properly these new places begin to feel unsafe and perceive the world around even with hostility, which led to a confinement in buildings they know and trust.
This confinement, which become like a prison, reduces a variety of chances, such as: physical exercise, the gathering of fresh food and the making of social activities.
A specific subset of people who experience this kind of problems every day are elders and, in order to freeing them of this burden, an indoor localization system for a wheeled walker is being researched . 

\section{Motivations}

There are many technical reasons to work in a project for assisted living, every one of them is a challenge and you can be considered a pioneer if your job is well-done.
However, personal satisfaction wasn't my primary goal when I had chosen this thesis and, whereas my colleagues felt the breeze derived from exploring a relatively new field, what I wanted to find was a solution to help real people who needs a "hand" from technology.
I realized that now it is the time to apply my effort in a product that is meaningful for a category of people who suffer. 
In fact, I could only imagine how hard is, for an old person who enjoyed his liberty when young, to be limited and have to rely on other people.
I can not say that I succeed because my work is just a small brick on the wall that this magnificent subject put in front of us, but I hope that my work would be a foundation for further work and that in a future not very far,a tool to assist elders will be invented .
For this, connected to the will of improving my skills, I undertook this path.

\section{The DALi project}

\vspace{1cm}
\begin{center}
      \includegraphics[width=0.3\textwidth]{img/Dali-logo.png}
\end{center}
\vspace{1cm}

The name stands for Devices for Assisted Living, which means the all the efforts of the team are focused in producing tools which helps people with deambulation or sensory problem in crowded places.
The project started in 2013 and it is composed by a consortium of universities and national research centers based in the following countries of Europe: United Kingdom, Spain, Greece, France and Austria.
Every partner contributes in specific areas of the project which, for example, involves both user centric design and human machine interfaces but also sensitive or cognitive algorithms.
The whole work is coordinated by The University of Trento and in particular by professor Luigi Palopoli.
At the moment of this writing, DALi is researching a new robotic wheeled walker, called rollator, which can help both the psychological sphere, by providing freedom of movement in newly discovered places to the user, and the physical one by actually aiding him during real movements.


\section{Role of the thesis inside DALi}
Designing and building a rollator brings many interesting challenges in the field of robotics, and even in the human computer interaction one, but they may reveal very difficult to achieve. 
In Particular, at the core of the walker functionalities there is a localization system which can track the user's position and provide directions to the user.
The precision requirement of the above mentioned unit is an extremely delicate matter and needed further investigation in order to bring the best possibilities whilst limiting the error.
Therefore the work described on the next pages explores the possibility of building an environment where qr-codes are the main providers of information about location and trajectory of the user.



\section{Outline}

1. Presentazione generale del progetto DALi e degli obbiettivi della tesi al suo interno (problema che la tesi risolve per il progetto)

2. descrizione del problema
        - marker nell’ambiente
        - necessità di riconoscerli e posizionarli rispetto alla telecamera
   stato dell’arte

3. Strumenti utilizzati (algoritmi di base, librerie)

4. descrizione dell’algoritmo

5. prove sperimentali di accuratezza in vari condizioni (su sistema per conto suo)

6. esperimenti sul carrelo (quando l’algoritmo è usato per la localizzazione)

7. Conclusioni






  
